\documentclass[cn,chinese,founder]{elegantbook}
\usepackage{wrapfig}
\everymath{\displaystyle}

\begin{document}
    \chapter{实数系的基本定理}
        \section{Cauthy收敛准则}
        
        \begin{definition}{基本数列}{Cauchy Sequence}
            称数列$\{x_n\}$为基本数列(或Cauthy数列),如果对每个$\varepsilon>0$,存在$N$,使得对每一对正整数$n,m>N$,成立估计式$\lvert a_n-a_m\rvert<\varepsilon$.
        \end{definition}
        
        \begin{theorem}{Cauthy收敛准则}{Cauthy Theorem} 
            数列收敛的充分必要条件是该数列为基本数列.
        \end{theorem}
        
        \begin{definition}{压缩映射}{Contraction Mapping}
            设函数$f$在区间$[a,b]$上定义,$f([a,b])\subset [a,b]$,并存在一个常数$k$,满足$0<k<1$,使得对一切 $x,y\in [a,b]$成立不等式$\lvert f(x)-f(y)\rvert\leqslant k\lvert x-y\rvert$,则称$f$是$[a,b]$上的一个压缩映射,称常数$k$为压缩常数.
        \end{definition}
        \begin{theorem}{压缩映射原理}{Contraction Mapping Theorem}
            设$f$是$[a,b]$上的一个压缩映射,则
            \begin{enumerate}
                \item $f$在$[a,b]$中存在唯一的不动点$\xi=f(\xi)$;
                \item 由任何初始值$a_0\in[a,b]$和递推公式$a_{n+1}=f(a_n),n\in N_+$生成的数列$\{a_n\}$一定收敛于$\xi$;
                \item 成立估计式$\lvert a_n-\xi\rvert\leqslant\frac{k}{1-k}\lvert a_n-a{n-1}\rvert$和$\lvert a_n-\xi\rvert\leqslant\frac{k^n}{1-k}\lvert a_1-a_0\rvert$(即事后估计与先验估计).
            \end{enumerate}
        \end{theorem}
            \subsection{思考题}
            \begin{example}
                Cauthy收敛准则在有理数集$\mathbb{Q}$中不不成立.
            \end{example}
            \begin{solution}
                原因和上一节相同,都是可能存在在实数系中的极限但是不在有理数集中.

                这个\href{https://www.zhihu.com/question/50995932/answer/866173110}{链接}有更加专业的解释,可以参考下图.

                \includegraphics[width=0.8\linewidth]{Cauthy收敛准则在Q中不成立.png}
            \end{solution}
        
            \subsection{练习题}
            \begin{exercise}
                满足以下条件的数列$\{x_n\}$是否一定是基本数列?若回答"是",请做出证明;若回答"不一定是",请举出反例:
                \begin{enumerate}
                    \item 对每个$\varepsilon>0$,存在$N$,当$n>N$是,成立
                    $\lvert x_n-x_N\rvert<\varepsilon$;
                    \begin{solution}
                        是.对于任意$\varepsilon>0$,由题设知可以得到$N$,使得当$n>N$时成立$\lvert x_n-x_N\rvert<\frac{\varepsilon}{2},\lvert x_m-x_N\rvert<\frac{\varepsilon}{2}$,则$\lvert x_n-x_m\rvert<\lvert x_n-x_N\rvert+\lvert x_m-x_N\rvert<\varepsilon$.
                    \end{solution}
                    \item 对所有$n,p\in N_+$,成立不等式$\lvert x_{n+p}-x_n\rvert\leqslant\frac{p}{n}$;
                    \begin{solution}
                        不一定是.显然基本数列满足这个条件,但例如$x_n=\sum_{i=1}^{n}\frac{1}{i}$,显然成立题设条件,甚至我们可以直接得到估计,但是我们知道$\lim_{n\to\infty}x_n=+\infty$,所以这不一定是基本数列.
                    \end{solution}
                    \item 对所有$n,p\in N_+$,成立不等式$\lvert x_{n+p}-x_n\rvert\leqslant\frac{p}{n^2}$;
                    \begin{solution}
                        是.根据题设,我们可以得到更精确的估计:$\lvert x_{n+p}-x_n\rvert<\lvert x_{n+p}-x_{n+p-1}\rvert+\cdots+\lvert x_{n+1}-x_{n}\rvert\leqslant\sum_{i=n}^{n+p-1}\frac{1}{i^2}$.根据我们之前得到的结论,$S_n=\sum_{i=1}^{\infty}\frac{1}{i^2}$收敛,所以只需$N$足够大就可以得到$\lvert x_{n+p}-x_n\rvert<S_{n+p}-S_n<\varepsilon$.
                    \end{solution}
                    \item 对每个正整数$p$,成立$\lim_{n\to\infty}(x_n-x_{n+p})=0$.
                    \begin{solution}
                        不一定是.例如(2)中的例子$x_n=\sum_{i=1}^{n}\frac{1}{i}$,对每个固定的$p$都满足题设,但它不是基本数列.
                    \end{solution}
                \end{enumerate}
            \end{exercise}

            \begin{exercise}
                用对偶法则于数列收敛的Cauthy收敛准则,以正面方式写出数列发散的充分必要条件.
            \end{exercise}
            \begin{solution}

                Cauthy收敛准则:$\forall \varepsilon>0,\exists N\in N_+$,当$ n,m>N\rightarrow\lvert x_n-x_m\rvert<\varepsilon\Longleftrightarrow \exists x\in R,\lim_{n\to\infty}x_n=x$;

                数列发散充要条件:数列$\{x_n\}$发散$\Longleftrightarrow \exists\varepsilon>0,\forall N\in N_+,\exists n,m>N,\lvert x_n-x_m\rvert\geqslant \varepsilon$.
            \end{solution}

            \begin{exercise}
                证明下列数列为基本数列,因此都是收敛数列:
                \begin{enumerate}
                    \item $a_n=1+\frac{1}{2!}+\frac{1}{3!}+\cdots+\frac{1}{n!},n\in N_+$;
                    \begin{solution}
                        不妨设$m\geqslant n>0$,而显然$\lvert a_m-a_n\rvert=\frac{1}{(n+1)!}+\cdots+\frac{1}{m!}<\frac{1}{2^{n+1}}+\cdots+\frac{1}{2^m}$,而后者是收敛的,知对任意$\varepsilon$都可以找到题目需要的$N$.
                    \end{solution}
                    \item $b_n=1-\frac{1}{2}+\frac{1}{3}-\cdots+(-1)^{n+1}\frac{1}{n},n\in N_+$;
                    \begin{solution}
                        $\lvert x_m-x_n\rvert=\left|(-1)^{n+2}\frac{1}{n+1}+\cdots+(-1)^{m+1}\frac{1}{m}\right|$,若$m-n$为奇数,则知从第二项开始可以每两项凑成一对,而且他们都与第一项正负相反;同时$\lvert x_m-x_n\rvert$与第一项同正负(因为最后一项和第一项同正负,而去掉第一项之后又可以两两凑成一对,这样的对子和第一项同正负),即$\big\lvert (-1)^{n+2}\frac{1}{n+1}+\cdots+(-1)^{m+1}\frac{1}{m}\big\rvert<\lvert \frac{1}{n+1}\rvert$.

                        例如$\left|1-\frac{1}{2}+\frac{1}{3}\right|=\underbrace{\left|(1-\frac{1}{2})+\frac{1}{3}\right|}_{\text{可以看出此式与1同号}}=\underbrace{\left|1+(-\frac{1}{2}+\frac{1}{3})\right|}_{\text{可以看出绝对值小于1}}$.

                        如果$m-n$为偶数,则两两凑成一对即知式子绝对值小于第一项绝对值.综上只需取$N$使得$\frac{1}{n+1}<\varepsilon$即可.
                    \end{solution}
                    \item $c_n=\frac{\sin2x}{2(2+\sin2x)}+\frac{\sin3x}{3(3+\sin3x)}+\cdots+\frac{\sin nx}{n(n+\sin nx)},n\in N_+$.
                    \begin{solution}
                        显然有$\left|x_m-x_n\right|\leqslant\left|\frac{1}{(n+1)(n)}+\frac{1}{m(m-1)}\right|$,裂项相消即可.
                    \end{solution}
                \end{enumerate}
            \end{exercise}

            \begin{exercise}
                设$a_n=\sin1+\frac{\sin2}{2!}+\cdots+\frac{\sin n}{n!},n\in N_+$,证明:
                \begin{enumerate}
                    \item 数列$\{a_n\}$有界,但不单调;
                    \begin{solution}
                        显然$a_n$不单调;而$a_n\leqslant\left|\sin1\right|+\cdots+\left|\frac{\sin n}{n!}\right|\leqslant\frac{1}{1!}+\cdots+\frac{1}{n!}$,由$\sum_{i=1}^\infty\frac{1}{i!}$收敛知$\{a_n\}$有界.
                    \end{solution}
                    \item $\{a_n\}$收敛.
                    \begin{solution}
                        任意两项之差绝对值可以按照(1)放缩,由$\sum_{i=1}^\infty\frac{1}{i!}$收敛知使差绝对值$<\varepsilon$的$N$是存在的,按照Cauthy收敛准则知$\{a_n\}$收敛.
                    \end{solution}
                \end{enumerate}
            \end{exercise}

            \begin{exercise}
                设从某个数列$\{x_n\}$定义$x_n=\sum_{k=1}^na_k,y_n=\sum_{k=1}^n\left|a_k\right|,n\in N_+$,若数列$\{y_n\}$收敛,证明数列$\{x_n\}$也收敛.

                (本题可以看成是上一题和例题3.4.1的推广.)
            \end{exercise}
            \begin{solution}
                有$\left|x_m-x_n\right|\leqslant y_m-y_n$,而因为$\{y_n\}$收敛,由Cauthy收敛准则知对任意$\varepsilon$可以找到$N$使$y_m-y_n<\varepsilon$,而根据前边不等式知这个$N$对$\{x_n\}$也成立,即$\{x_n\}$也是基本数列,故$\{x_n\}$收敛.
            \end{solution}

            \begin{exercise}
                设$S_n=1+\frac{1}{2^p}+\frac{1}{3^p}+\cdots+\frac{1}{n^p},n\in N_+$,其中$p\leqslant 1$,证明$\{S_n\}$发散.
            \end{exercise}
            \begin{solution}
                
            \end{solution}
\end{document}